\documentclass[conference]{IEEEtran}
\IEEEoverridecommandlockouts

\usepackage{cite}
\usepackage{amsmath,amssymb,amsfonts}
\usepackage{graphicx}
\usepackage{textcomp}
\usepackage{xcolor}
\usepackage{listings}
\usepackage{url}
\usepackage{booktabs}
\usepackage{hyperref}

\lstset{
  basicstyle=\ttfamily\footnotesize,
  breaklines=true,
  breakatwhitespace=true,
  columns=flexible,
  keepspaces=true,
  frame=single,
  framerule=0pt,
  backgroundcolor=\color{black!4},
  xleftmargin=4pt,
  xrightmargin=4pt,
  aboveskip=4pt,
  belowskip=4pt,
}

\begin{document}

\title{Fingerprinting and Obfuscating WireGuard:\\
A Practical Demonstration of Protocol Identification and udp2raw Evasion}

\author{\IEEEauthorblockN{Steven Schiavone}
\IEEEauthorblockA{\textit{Seidenberg School of CSIS} \\
\textit{Pace University}\\
New York, NY \\
sschiavone@pace.edu}
\and
\IEEEauthorblockN{Nicholas Johannsen}
\IEEEauthorblockA{\textit{Seidenberg School of CSIS} \\
\textit{Pace University}\\
New York, NY}
}

\maketitle

\begin{abstract}
WireGuard is widely deployed for its simplicity, performance, and modern
cryptography. We show, however, that its over-the-wire framing is trivially
identifiable: a constant four-byte header structure and three fixed-length
handshake packets allow a passive observer to fingerprint a WireGuard session
without decrypting a single byte. We deploy a WireGuard server on Oracle Cloud
Infrastructure (OCI), connect a Kali Linux client through it, and use a Scapy
sniffer of our own design to identify 102 WireGuard packets in a 25-second
capture, including the handshake initiation and response. We then wrap the
same tunnel in \textit{udp2raw} \texttt{faketcp} obfuscation. The tunnel
continues to function (5/5 ICMP pings, 0\% loss to the inner peer) but the
identical Scapy sniffer reports zero WireGuard packets. We conclude that
WireGuard's straightforward design --- a feature in trusted networks ---
becomes a liability under deep-packet-inspection (DPI) regimes, and that a
small obfuscation layer is sufficient to defeat the trivial fingerprint.
\end{abstract}

\begin{IEEEkeywords}
WireGuard, VPN, deep packet inspection, traffic obfuscation, udp2raw,
Scapy, network forensics
\end{IEEEkeywords}

\section{Introduction}

A central tension of modern VPN design is that strong end-to-end cryptography
does not, by itself, conceal the presence of a VPN. Censorship and
surveillance regimes from corporate firewalls to nation-state DPI middleboxes
typically do not attempt to break the encryption; they fingerprint the
\textit{protocol} and drop or throttle the flow. WireGuard~\cite{donenfeld2017},
prized in the research community for its small attack surface and modern
constructions, has a particularly stark version of this problem: its packets
follow a fixed and predictable layout that needs no machine learning to
identify.

This paper demonstrates the problem and a practical countermeasure. We deploy
a complete lab environment, build a minimal Python/Scapy classifier to
fingerprint WireGuard handshakes by their byte signature, and then evaluate
the same classifier against a tunnel obfuscated with
\textit{udp2raw}~\cite{wangyu2017udp2raw} in \texttt{faketcp} mode. The
contribution is not novel cryptography but a clean reproduction of the gap
between ``encrypted'' and ``unidentifiable.''

\section{Background}

\subsection{The WireGuard wire format}

The WireGuard whitepaper~\cite{donenfeld2017} specifies four message types,
each carried in a single UDP datagram. Every message begins with a one-byte
type identifier and three reserved zero bytes:

\begin{lstlisting}
+---------+--------------------+
| type(1) | reserved zero (3) |  ...
+---------+--------------------+
\end{lstlisting}

The handshake messages are also \textit{fixed length}, summarized in
Table~\ref{tab:wgtypes}.

\begin{table}[h]
\centering
\caption{WireGuard UDP payload signature.}
\label{tab:wgtypes}
\begin{tabular}{@{}lll@{}}
\toprule
Type byte & Message & UDP payload length \\
\midrule
0x01 & Handshake Initiation & 148 bytes \\
0x02 & Handshake Response   &  92 bytes \\
0x03 & Cookie Reply         &  64 bytes \\
0x04 & Transport Data       & variable, $\geq 32$ bytes \\
\bottomrule
\end{tabular}
\end{table}

The combination is unusually distinctive. Three concatenated checks ---
``UDP datagram, first byte is one of $\{1,2,3\}$, next three bytes are zero,
length is exactly 148/92/64'' --- yield a handshake fingerprint with negligible
false-positive risk in normal traffic. No state, no machine learning, no
heuristics; a 12-line classifier suffices.

\subsection{udp2raw and traffic obfuscation}

\textit{udp2raw}~\cite{wangyu2017udp2raw} is a single-binary tunnel that
takes UDP datagrams on one side and re-emits them as crafted raw packets that
imitate other protocols (TCP, ICMP, or fake-UDP). In \texttt{faketcp} mode,
udp2raw performs a SYN/ACK handshake with a corresponding instance on the
other side, then exchanges encrypted payloads inside what appears to be an
ordinary long-running TCP session. The wrapped UDP datagrams --- including
the WireGuard handshake messages we wish to hide --- never appear on the
wire as UDP at all.

We chose \texttt{faketcp} over the alternative modes because it produces the
strongest disguise (it survives more middleboxes than ICMP or fake-UDP) and
is widely deployed in the censorship-circumvention community.

\section{Lab environment}

The experimental setup, illustrated conceptually in Fig.~\ref{fig:topo},
consists of three nodes:

\begin{enumerate}
  \item A WireGuard server on an OCI VM.\@ Oracle Linux 9.7 on aarch64
        (VM.Standard.A1.Flex), public IP 158.101.122.42. WireGuard listens
        on UDP/51820; for the obfuscated path, udp2raw listens on TCP/4096
        and forwards locally to UDP/51820.
  \item A Kali Linux 2025.3 client running in VirtualBox on a Windows 11
        host. Network adapter set to NAT (10.0.2.15) since Wi-Fi bridging
        is unreliable across most consumer Wi-Fi NICs. Internet egress is
        through the Windows host.
  \item A Scapy-based DPI tool running inside the Kali VM, sniffing the
        VM's NAT interface. This is the ``passive observer'' position
        from a censorship-modelling perspective.
\end{enumerate}

\begin{figure}[h]
\centering
\fbox{\parbox{0.95\linewidth}{\footnotesize\ttfamily
\textbf{Direct path (DPI-visible):}\\
~Kali wg client $\xrightarrow{\text{UDP/51820}}$ OCI wg server\\[3pt]
\textbf{Obfuscated path (DPI-blind):}\\
~Kali wg client $\xrightarrow{\text{lo:51821}}$ Kali udp2raw client\\
~$\xrightarrow{\text{TCP/4096 (faketcp)}}$ OCI udp2raw server\\
~$\xrightarrow{\text{lo:51820}}$ OCI wg server
}}
\caption{Two end-to-end paths under test. The wireguard interface and keys
are identical in both cases; only the transport between client and server
changes.}
\label{fig:topo}
\end{figure}

The client maintains two WireGuard configurations (\texttt{wg-direct} and
\texttt{wg-obfuscated}) using identical keys and identical inner addressing
(10.66.66.2/24); only the \texttt{Endpoint} differs. This isolates the
obfuscation as the sole independent variable.

\subsection{Path MTU considerations}

Stacking WireGuard inside udp2raw inside Ethernet introduces a non-trivial
MTU constraint. Each layer subtracts overhead: WireGuard adds 80 bytes for
its header, AEAD nonce, and Poly1305 tag; udp2raw \texttt{faketcp} adds
roughly 60 bytes of TCP header and obfuscation framing. The default
WireGuard MTU of 1420 bytes is therefore 80--100 bytes too large for the
obfuscated path. Symptomatically, packets just under the link MTU are
silently dropped by udp2raw with a \texttt{huge packet, data
len=1472 (>=1375)} warning, and the faketcp ``connection'' is torn down by
intermediate conntrack timeouts roughly every 17 seconds.

We resolved this by setting \texttt{MTU = 1200} on the obfuscated WireGuard
interface, well within udp2raw's effective payload budget. End-to-end
ICMP through the tunnel then runs at $\sim$14\,ms RTT with 0\% loss across
30-packet bursts.

\section{The Scapy DPI classifier}

Our classifier consists of two short Python files: a stateless predicate that
maps a UDP payload to one of the four WireGuard message classes, and a
sniffer that applies that predicate to live or recorded traffic. The complete
predicate is reproduced in Listing~\ref{lst:classify}.

\begin{lstlisting}[caption={The complete WireGuard fingerprint predicate, less than 25 lines.},label={lst:classify}]
WG_TYPES = {
    1: ("HANDSHAKE_INIT",     148),
    2: ("HANDSHAKE_RESPONSE",  92),
    3: ("COOKIE_REPLY",        64),
    4: ("TRANSPORT_DATA",    None),
}

def classify(payload):
    if len(payload) < 4:
        return None
    msg_type = payload[0]
    if msg_type not in WG_TYPES \
       or payload[1:4] != b"\x00\x00\x00":
        return None
    label, expected_len = WG_TYPES[msg_type]
    if expected_len is not None:
        if len(payload) != expected_len:
            return None
        return label, msg_type, True
    if len(payload) < 24:
        return None
    return label, msg_type, False
\end{lstlisting}

The third tuple element distinguishes \textit{exact} matches (handshake
packets that pass the strict length check) from \textit{shape} matches
(transport-data packets that pass only the type+reserved check). To
sanity-check the predicate, we wrote a self-test of eight synthetic
payloads --- four valid messages of each type and four near-misses (truncated,
wrong type byte, non-zero reserved bytes, and a type-1 message of incorrect
length). All eight cases pass.

The sniffer wraps Scapy's \texttt{sniff()} (or \texttt{rdpcap()} for offline
analysis), filters to UDP, applies the predicate, and produces both a
per-packet log and an aggregate summary. The orchestrator script
\texttt{run\_demo.sh} captures a pcap with \texttt{tcpdump}, generates traffic
through whichever WireGuard interface is currently up, and replays the pcap
through the sniffer for the canonical evidence file.

\section{Results}

\subsection{Direct WireGuard: easily identified}

With the \texttt{wg-direct} configuration active, we ran the orchestrator
with a 25-second capture window. The handshake completed in 1 second,
five-of-five pings succeeded, and the curl through the tunnel reported
the correct origin IP (158.101.122.42).

The DPI summary output from \texttt{wg\_dpi.py} on the captured pcap is
reproduced in Listing~\ref{lst:baseline}.

\begin{lstlisting}[caption={DPI run on the direct WireGuard path.},label={lst:baseline}]
=== Summary ===
Total WireGuard packets detected: 102
  TRANSPORT_DATA       100
  HANDSHAKE_INIT       1
  HANDSHAKE_RESPONSE   1

Top flows:
  10.0.2.15:35004 -> 158.101.122.42:51820   58
  158.101.122.42:51820 -> 10.0.2.15:35004   43
\end{lstlisting}

The handshake is unambiguously located. Per-packet output (omitted for space)
shows the 148-byte initiation, the 92-byte response, the typical sequence of
small (96--128~byte) keepalive-class transport packets, and bursts of
1344--1456~byte packets corresponding to bulk traffic such as the curl
response. None of this requires any state across packets.

\subsection{udp2raw: indistinguishable from background TCP}

With \texttt{wg-direct} torn down and \texttt{wg-obfuscated} brought up, we
re-ran the same orchestrator with the BPF filter widened to capture all
traffic on TCP/4096 (the udp2raw transport port). The faketcp handshake took
11 seconds to settle; once established, ping again succeeded
five-of-five at 0\% loss. The DPI summary is reproduced in
Listing~\ref{lst:obfuscated}.

\begin{lstlisting}[caption={DPI run on the udp2raw-obfuscated path.},label={lst:obfuscated}]
=== Summary ===
No WireGuard traffic detected.
\end{lstlisting}

The pcap is not empty: it contains tens of TCP segments on port 4096 with
plausible sequence numbers and ACK behaviour. None of those segments contain
a UDP payload starting \texttt{0x01 0x00 0x00 0x00}, none are 148/92/64
bytes, and the WireGuard handshake messages are nowhere on the wire.

The Scapy classifier, exactly as written in
Listing~\ref{lst:classify}, reports zero matches. The functional VPN tunnel
is intact end-to-end.

\subsection{Comparison}

Table~\ref{tab:results} contrasts the two runs. The defining observation is
that all functional metrics --- ping success, end-to-end RTT --- are
preserved or barely changed, while the DPI metric collapses entirely.

\begin{table}[h]
\centering
\caption{End-to-end behaviour and observability of the two paths.}
\label{tab:results}
\begin{tabular}{@{}lcc@{}}
\toprule
Metric & Direct & udp2raw \\
\midrule
WireGuard packets identified                & 102 & 0 \\
Handshake packets (INIT + RESPONSE)         & 2   & 0 \\
Inner-tunnel ICMP success                   & 5/5 & 5/5 \\
Inner-tunnel mean RTT                       & 33.9~ms & 23.1~ms \\
Time to handshake                           & $<$1\,s & 11\,s \\
Wire MTU on the wg interface                & 1420 & 1200 \\
\bottomrule
\end{tabular}
\end{table}

\section{Discussion}

\subsection{What is and isn't shown}

The fingerprint we exploit is not a vulnerability in WireGuard's
cryptography. The session keys, the AEAD construction, and the noise-protocol
handshake remain intact under both configurations. We exploit only that
WireGuard does not attempt to be \textit{indistinguishable} from other
traffic. This is by design: WireGuard's whitepaper~\cite{donenfeld2017}
explicitly states that obfuscation against censorship is out of scope and
should be handled by an external layer if needed.

It is worth noting that the udp2raw obfuscation defeats only the
\textit{passive} fingerprint we constructed. An adversary controlling
endpoints between the udp2raw client and server could perform active probing
(e.g., sending crafted SYN packets to TCP/4096 and observing replies) and
likely identify the udp2raw shim itself. Defeating active probing requires
domain-fronting or similar protocol-mimicry tools beyond udp2raw.

\subsection{Implications}

For the typical home or office WireGuard deployment --- where the goal is
encryption and not concealment --- our result simply confirms an existing
design property. For users in restrictive networks, the result has more
weight: the same property that makes WireGuard easy to audit makes it easy
to block. The widely-cited Great Firewall of China is documented to identify
and block WireGuard handshakes, and other regimes are presumed to do the
same. A 25-line Python script, one minute of capture, and one page of
analysis yield enough signal to drop every WireGuard tunnel a national-scale
DPI box might encounter --- absent obfuscation.

\subsection{Reproducibility}

All scripts, configurations, and evidence files used to produce this paper
are available in the project repository. The single shell command
\texttt{run\_demo.sh} reproduces the headline results in approximately
50 seconds end-to-end on a clean Kali install. The two evidence files
discussed above (\texttt{20260508-203817\_baseline-pcap.txt} and
\texttt{20260508-203845\_udp2raw-pcap.txt}) are timestamped artifacts of the
runs reported here.

\section{Conclusion}

WireGuard provides excellent confidentiality and integrity, but no
unobservability. A passive observer with no decryption capability and a
classifier of fewer than 25 lines reliably identifies WireGuard sessions in
practice. Wrapping the tunnel in a small obfuscation layer ---
udp2raw \texttt{faketcp} in our experiments --- removes that fingerprint
without measurably degrading the user experience, at a small cost in MTU
and a one-time handshake-establishment delay.

For VPN users in restrictive network environments, a kernel-mode WireGuard
tunnel without an obfuscation shim is, in effect, broadcasting that it is
a VPN.\@ For VPN users who only require confidentiality on a trusted network,
the unwrapped tunnel is sufficient and faster. The choice between the two is
a deployment decision, and we hope this paper has made the underlying
trade-off concrete.

\begin{thebibliography}{00}
\bibitem{donenfeld2017} J. A. Donenfeld, ``WireGuard: Next Generation Kernel
Network Tunnel,'' \textit{NDSS Symposium 2017}, San Diego, CA, USA, 2017.

\bibitem{wangyu2017udp2raw} ``udp2raw-tunnel: A tunnel which turns UDP
traffic into encrypted UDP/FakeTCP/ICMP traffic,'' wangyu-, GitHub
repository, \url{https://github.com/wangyu-/udp2raw}, accessed 2026-05-08.

\bibitem{scapy} P. Biondi and the Scapy community, ``Scapy: Packet
crafting for Python2 and Python3,'' \url{https://scapy.net}.
\end{thebibliography}

\end{document}
